\documentclass{beamer}
\usepackage{bussproofs}


%------------------------------------------------
%   MOJE DEFINICIJE
%------------------------------------------------

\newcommand{\type}{\text{ type}}
\newcommand{\mc}[1]{\mathcal{#1}}
\newcommand{\cd}[1]{\texttt{#1}}
\newcommand{\G}{\Gamma}
\newcommand{\D}{\Delta}
\newcommand{\os}[2]{\onslide<#1>{#2}}

\newcommand{\AXC}[1]{\AxiomC{#1}}
\newcommand{\UIC}[1]{\UnaryInfC{#1}}
\newcommand{\BIC}[1]{\BinaryInfC{#1}}
\newcommand{\TIC}[1]{\TrinaryInfC{#1}}
\newcommand{\RL}[1]{\RightLabel{#1}}
\newcommand{\DP}{\DisplayProof}


%------------------------------------------------
%   PACKAGES
%------------------------------------------------
\usepackage{xcolor}
\usepackage{tikz}
\usetikzlibrary{shadings}
\usepackage[most]{tcolorbox}

% \EnableBpAbbreviations

\usefonttheme{professionalfonts}


%------------------------------------------------
%   COLORS & TEXT
%------------------------------------------------
\definecolor{myviolet}{RGB}{240, 121, 242}
% \definecolor{lightviolet}{RGB}{255, 200, 255}
\definecolor{lightviolet}{RGB}{255, 179, 255}
\definecolor{lighterviolet}{RGB}{252, 220, 252}


\setbeamercolor{normal text}{fg=black,bg=myviolet}

\setbeamertemplate{navigation symbols}{}

%------------------------------------------------
%   GLAVNI NASLOV
%------------------------------------------------

% Rainbow-bordered title box (reuses rainbow math border)
\tcbset{
    rainbow title/.style={
        enhanced,
        colback=lightviolet,     % background inside the box
        coltext=black,
        arc=6pt,
        outer arc=6pt,
        boxrule=2pt,
        opacityback=1,
        frame code={
            % draw the rainbow border (same as rainbow math)
            \shade[shading=rainbowshading, rounded corners=6pt]
                (frame.south west) rectangle (frame.north east);
        },
        interior code={
            % draw inner light-violet box
            \fill[lightviolet, rounded corners=4pt]
                (interior.south west) rectangle (interior.north east);
        },
        left=12pt, right=12pt, top=10pt, bottom=10pt,
    }
}

% Title page with rainbow-bordered box
\defbeamertemplate*{title page}{rainbowstyle}{
    \begin{center}
        \begin{tcolorbox}[rainbow title]
            \centering
            {\ttfamily\bfseries\Huge\inserttitle}
        \end{tcolorbox}

        \vspace{1em}

        {\large\insertauthor}

        \vspace{0.5em}

        {\small\insertdate}
    \end{center}
}



% \setbeamercolor{frametitle}{bg=lightviolet, fg=black}
% \setbeamertemplate{frametitle}{
%     \begin{beamercolorbox}[ht=1.8em,wd=\paperwidth]{frametitle}
%         \centering
%         {\ttfamily\bfseries\Large\insertframetitle}
%     \end{beamercolorbox}
% }
%------------------------------------------------
%   FRAMETITLE
%------------------------------------------------

\usepackage{tikz}
\usepackage[most]{tcolorbox}

\makeatletter
\tcbset{
    rainbow frametitle/.style={
        % enhanced,
        colback=lightviolet,
        coltext=black,
        % outer arc=0pt,
        boxrule=0pt,
        opacityback=1,
        left=0pt, right=0pt, top=10pt, bottom=10pt,
        width=\paperwidth-8pt,
        frame code={
        },
        overlay={
            % position the box absolutely at the top of the page
            \path[shift={(current page.north west)}] 
                (0,0) rectangle (\paperwidth,-\ht\strutbox);
        }        
    }
}

\setbeamertemplate{frametitle}{
    % use a tcolorbox absolutely positioned
    \begin{tikzpicture}[remember picture,overlay]
        \node[anchor=north west, xshift=-1pt, yshift=1pt] at (current page.north west) {
            \begin{tcolorbox}[rainbow frametitle]
                \centering
                {\ttfamily\bfseries\Large\insertframetitle}
            \end{tcolorbox}
        };
    \end{tikzpicture}
}
\makeatother





% \setbeamertemplate{frametitle}{width={\dimexpr\paperwidth-2\slideinset\relax}}


% \setbeamertemplate{frametitle}{%
%     \noindent\hspace*{\slideinset}% place at inner left edge (do not cover outer border)
%     \begin{tcolorbox}[rainbow frametitle,width={\dimexpr\paperwidth-2\slideinset\relax}]
%         	tfamily\bfseries\large\insertframetitle
%     \end{tcolorbox}
% }





%------------------------------------------------
%   RAINBOW SHADING
%------------------------------------------------
\pgfdeclarehorizontalshading{rainbowshading}{100bp}{
	color(0bp)=(red);
    color(20bp)=(orange);
    color(40bp)=(yellow);
    color(60bp)=(green);
    color(80bp)=(blue);
    color(100bp)=(violet)
}

%------------------------------------------------
%   SLIDE BACKGROUND (RAINBOW BORDER)
%------------------------------------------------
\usebackgroundtemplate{
\begin{tikzpicture}[remember picture,overlay]

    \fill[myviolet]
        (current page.south west) rectangle (current page.north east);

    \shade[shading=rainbowshading]
        (current page.south west) rectangle (current page.north east);

    \fill[myviolet]
        ([xshift=4pt,yshift=4pt]current page.south west)
        rectangle
        ([xshift=-4pt,yshift=-4pt]current page.north east);

\end{tikzpicture}
}



%------------------------------------------------
%   RAINBOW-BORDERED MATH ENVIRONMENT
%------------------------------------------------
\tcbset{
    rainbow math/.style={
        enhanced,
        colback=white,
        % colback=lightviolet,
        coltext=black,
        arc=6pt,                            % rounded corners
        outer arc=6pt,
        boxrule=2pt,
        opacityback=1,
        frame code={
            % draw the rainbow border
            \shade[shading=rainbowshading, rounded corners=6pt]
                (frame.south west) rectangle (frame.north east);
        },
        interior code={
            % draw inner white box
            \fill[lighterviolet, rounded corners=4pt]
                (interior.south west) rectangle (interior.north east);
        },
        left=10pt, right=10pt, top=8pt, bottom=8pt,
    }
}

\newenvironment{whitemath}
{\begin{tcolorbox}[rainbow math]}
{\end{tcolorbox}}




\title{Linearni tipi}
\author{Vida Mlinar}
\institute[FMF]{Fakulteta za matematiko in fiziko}
\date{}


%------------------------------------------------
%   DOCUMENT
%------------------------------------------------
\begin{document}

\frame{\maketitle}

\begin{frame}{Klasična logika}
    \begin{whitemath}
        \begin{center}
            \AXC{\(\G \vdash P\)}
            \AXC{\(\G \vdash Q\)}
            \RL{\(\land\)-intro}
            \BIC{\(\G \vdash P \land Q\)}
            \DP\\
            \os{2-}{\AXC{\(\G \vdash P \land Q\)}
            \RL{\(\land_1\)-elim}
            \UIC{\(\G \vdash P\)}
            \DP}
            \quad
            \os{3-}{\AXC{\(G \vdash P \land Q\)}
            \RL{\(\land_2\)-elim}
            \UIC{\(\G \vdash Q\)}
            \DP}
        \end{center}
    \end{whitemath}
    \os{4-}{\begin{whitemath}
        \begin{center}
            \os{6-}{\AXC{\(\G, P \vdash Q\)}
            \RL{\(\Rightarrow\)-intro}
            \UIC{\(\G \vdash P \Rightarrow Q\)}
            \DP}
            \quad
            \os{5-}{\AXC{\(\G \vdash P \Rightarrow Q\)}
            \AXC{\(\G \vdash P\)}
            \RL{\(\Rightarrow\)-elim}
            \BIC{\(\G \vdash Q\)}
            \DP}
        \end{center}
    \end{whitemath}}
    
    \begin{whitemath}
        \begin{prooftree}
            \AXC{\(P \in \G\)}
            \UIC{\(\G \vdash P\)}
        \end{prooftree}
    \end{whitemath}
\end{frame}


\begin{frame}
    \begin{whitemath}
        \begin{center}
            \AXC{}
            \UIC{\(P \in P\)}
            \UIC{\(P \vdash P\)}
            \AXC{}
            \UIC{\(P \in P\)}
            \UIC{\(P \vdash P\)}
            \BIC{\(P \vdash P \land P\)}
            \DP
            {\huge\(\rightsquigarrow\)}
            \AXC{}
            \UIC{\(P \vdash P \land P\)}
            \DP
        \end{center}
            
    
    \end{whitemath}
    
    \begin{whitemath}
        
        \begin{center}
            \AXC{}
            \UIC{\(A \in \G \cup A\)}
            \UIC{\(\G, A \vdash A\)}
            \AXC{\(\G \vdash B\)}
            \BIC{\(\G, A \vdash A \land B\)}
            \UIC{\(\G, A \vdash B\)}
            \DP
            {\huge\(\rightsquigarrow\)}
            \AXC{\(\G \vdash B\)}
            \UIC{\(\G, A \vdash B\)}
            \DP
        \end{center}
    \end{whitemath}
\end{frame}

% \begin{frame}{Linearna logika}
%     \begin{whitemath}
%         \begin{center}
%             \AXC{\(\G \vdash P\)}
%             \AXC{\(\D \vdash Q\)}
%             \RL{\(\land\)-intro}
%             \BIC{\(\G, \D \vdash P \land Q\)}
%             \DP\\
%             \AXC{\(\G \vdash P \land Q\)}
%             \RL{\(\land_1\)-elim}
%             \UIC{\(\G \vdash P\)}
%             \DP
%             \quad
%             \AXC{\(\G \vdash P \land Q\)}
%             \RL{\(\land_2\)-elim}
%             \UIC{\(\G \vdash Q\)}
%             \DP
%         \end{center}
%     \end{whitemath}
%     \begin{whitemath}
%         \begin{center}
%             \AXC{\(\G, P \vdash Q\)}
%             \RL{\(\Rightarrow\)-intro}
%             \UIC{\(\G \vdash P \Rightarrow Q\)}
%             \DP
%             \quad
%             \AXC{\(\G \vdash P \Rightarrow Q\)}
%             \AXC{\(\D \vdash P\)}
%             \RL{\(\Rightarrow\)-elim}
%             \BIC{\(\G \vdash Q\)}
%             \DP
%         \end{center}
%     \end{whitemath}
    
%     \begin{whitemath}
        
%         \begin{prooftree}
%             \AXC{\(P \in \G\)}
%             \UIC{\(\G \vdash P\)}
%         \end{prooftree}
%     \end{whitemath}
% \end{frame}

% \begin{frame}
%     Klasična logika:
%     \begin{whitemath}
%         \begin{prooftree}
%             \AXC{\(P \in P\)}
%             \UIC{\(P \vdash P\)}
%             \AXC{\(P \in P\)}
%             \UIC{\(P \vdash P\)}
%             \BIC{\(P \vdash P \land P\)}
%         \end{prooftree}
%     \end{whitemath}
%     %
%     Linearna logika:
%     \begin{whitemath}
%         \begin{prooftree}
%             \AXC{\(P \in P\)}
%             \UIC{\(P \vdash P\)}
%             \AXC{\(P \in P\)}
%             \UIC{\(P \vdash P\)}
%             \BIC{\(P, P \vdash P \land P\)}
%         \end{prooftree}
%     \end{whitemath}
% \end{frame}


\begin{frame}
    
    \begin{whitemath}
        
        \begin{prooftree}
            \AXC{\(x: A \in \G\)}
            \UIC{\(\G \vdash x:A\)}
        \end{prooftree}
    \end{whitemath}
    \begin{whitemath}
        \begin{center}
            \AXC{\(\G \vdash M_1 : A_1\)}
            \AXC{\(\G \vdash M_2 : A_2\)}
            \BIC{\(\G \vdash (M_1, M_2) : A_1 \times A_2\)}
            \DP\\\ \\
            \AXC{\(\G \vdash M: A_1 \times A_2\)}
            \UIC{\(\G \vdash fst(M): A_1\)}
            \DP
            \qquad
            \AXC{\(\G \vdash M: A_1 \times A_2\)}
            \UIC{\(\G \vdash snd(M): A_2\)}
            \DP
        \end{center}
    \end{whitemath}
    \begin{whitemath}
    \begin{center}
        \AXC{\(\G, x:A \vdash M:B\)}
        \UIC{\(\G \vdash (x \mapsto M): A \rightarrow B\)}
        \DP\\\ \\
        \AXC{\(\G M_1: A \rightarrow B\)}
        \AXC{\(\G \vdash M_2: A\)}
        \BIC{\(\G \vdash M_1 M_2 : B\)}
        \DP
    \end{center}
        
    \end{whitemath}
\end{frame}

% \section{Teorija tipov}
% 

\begin{frame}{Teorija tipov}

\begin{whitemath}
\begin{align*}
    \Gamma &\vdash A \type \\
    \Gamma &\vdash A \doteq B \type \\
    \Gamma &\vdash a : A \\
    \Gamma &\vdash a \doteq b: A
\end{align*}
\end{whitemath}

\end{frame}
\begin{frame}[plain]
    \begin{whitemath}
    \begin{prooftree}
        \AXC{\(A \type\)}
        \AXC{\(B \type\)}
        \RL{\(\rightarrow\) form}
        \BIC{\(A \rightarrow B\)}
    \end{prooftree}
    \begin{prooftree}
        \AXC{\(a:A\)}
        \AXC{\(b:B\)}
        \RL{\(\rightarrow \mathcal{I}\)}
        \BIC{\(\lambda x:A. b : A \rightarrow B\)}
    \end{prooftree}
    \begin{prooftree}
        \AXC{\(f: A \rightarrow B\)}
        \AXC{\(a:A\)}
        \RL{\(\rightarrow \mathcal{E}\)}
        \BIC{\(f~a : B\)}
    \end{prooftree}
    \end{whitemath}
    \onslide<2->{
    \begin{whitemath}
        \begin{prooftree}
            \AXC{\(\onslide<3>{\Gamma \vdash} A \type\)}
            \AXC{\(\onslide<3>{\Gamma \vdash} B \type\)}
            \RL{\(K\) form}
            \BIC{\(\onslide<3>{\Gamma \vdash} C_1\ A\ |\ C_2\ B \type\)}
        \end{prooftree}
        % \begin{prooftree}
            
            \begin{equation*}
                \AXC{\(\onslide<3>{\Gamma \vdash} a: A\)}
                \RL{\(K \mc{I}\)}
                \UIC{\(\onslide<3>{\Gamma \vdash} C\ a : K\)}
                \DP
                \qquad
                \AXC{\(\onslide<3>{\Gamma \vdash} b:B\)}
                \RL{\(K \mc{I}\)}
                \UIC{\(\onslide<3>{\Gamma \vdash} C\ b : K\)}
                \DP
            \end{equation*}
        % \end{prooftree}
        \begin{prooftree}
            \AXC{\(\onslide<3>{\Gamma \vdash} u:K\)}
            \AXC{\(\onslide<3>{\Gamma,}~a:A \vdash v:V\)}
            \AXC{\(\onslide<3>{\Gamma,}~b:B \vdash w:V\)}
            \RL{\(K\mc{E}\)}
            \TIC{\(\onslide<3>{\Gamma \vdash} (\cd{case}\ u\ \cd{of}\ C_1 a \rightarrow v\ |\ C_2\ b \rightarrow w) : V\)}
        \end{prooftree}}
    \end{whitemath}
\end{frame}


\begin{frame}{Primeri}
    \begin{whitemath}
        \begin{prooftree}
            \AXC{\(\cdot\)}
            \RL{\cd{Bool}\ \text{form}}
            \UIC{\cd{Bool} \type}
        \end{prooftree}
        \begin{align*}
            \AXC{\(\cdot\)}
            \RL{\(\cd{Bool}\ \mc{I}_1\)}
            \UIC{\cd{True} : \cd{Bool}}
            \DP
            \qquad
            \AXC{\(\cdot\)}
            \RL{\(\cd{Bool}\ \mc{I}_2\)}
            \UIC{\cd{False} : \cd{Bool}}
            \DP
        \end{align*}
    \end{whitemath}

    \onslide<2->{\begin{whitemath}
        
        \begin{prooftree}
            \AXC{\(A \type\)}
            \UIC{\(A\ \cd{List} \type\)}
        \end{prooftree}
        \begin{align*}
            \AXC{\(A \type\)}
            \UIC{\(\cd{Nil} : A\ \cd{List}\)}
            \DP
            \qquad
            \AXC{\(A \type\)}
            \UIC{\(\cd{Cons} : A \rightarrow A\ \cd{List} \rightarrow A\ \cd{List}\)}
            \DP
        \end{align*}
    \end{whitemath}}
\end{frame}


\begin{frame}{Pravilo šibenja \only<2>{in pravilo rezanja}}
    \begin{whitemath}
        \begin{align*}
            \AXC{\(\Gamma \vdash \Delta\)}
            \RL{\(\mc{LW}\)}
            \UIC{\(\Gamma,\ A \vdash \Delta\)}
            \DP
            \qquad
            \AXC{\(\Gamma \vdash \Delta\)}
            \RL{\(\mc{RW}\)}
            \UIC{\(\Gamma \vdash \Delta,\ A\)}
            \DP
        \end{align*}
    \end{whitemath}

    \onslide<2->{\begin{whitemath}
        \begin{align*}
            \AXC{\(\Gamma \vdash A\)}
            \AXC{\(\Delta, A \vdash B\)}
            \RL{cut}
            \BIC{\(\Gamma, \Delta \vdash B\)}
            \DP
        \end{align*}
    \end{whitemath}}
\end{frame}

\section{Linearni tipi}

\end{document}
